\documentclass[11pt]{article}
\usepackage[margin=1in]{geometry}
\usepackage{amsmath,amssymb,amsthm,hyperref}
\newtheorem{proposition}{Proposition}
\title{Long-Run Monopoly Selection and Sectoral Wages}
\author{Working derivations for Oliver Pardo}
\date{22 November 2026}
\begin{document}
\maketitle
I use the environment of \texttt{static\_equilibrium.tex}. This note treats
sector-specific workers with fixed positive shares $\ell_S,\ell_C$.
It proves a static selection result and conditional equilibrium implications,
not existence of the extended dynamic equilibrium.

\section{The global monopoly optimum at high capital intensity}
Let $H=AN$, $H_j=\ell_jH$, $k=K/H$, and $G(X;H)=H g(X/H)$.
Write $c=\omega_{XS}^{\sigma_S/(\sigma_S-1)}$ and
$\nu=(\sigma_S-1)/\sigma_S\in(0,1)$.
For a supply-to-capital ratio $y=X/K$, operating profit divided by $K$ is
\begin{equation}
 f_k(y;B)=(1-\alpha)
 \left[\frac{g(ky)}k\right]^{-\alpha}g'(ky)y-\frac yB.
 \label{eq:selection-objective}
\end{equation}
The value at $y=0$ is defined by continuous extension.

\begin{proposition}[Eventual uniqueness of global monopoly supply]
For $B$ in a sufficiently small compact neighborhood of any finite
$B_\infty>0$, the static monopoly optimum is unique whenever $k$ is
sufficiently large. Along any $k\to\infty$, $B\to B_\infty$,
\begin{equation}
 \frac{X^*}{K}\to y_*(B_\infty)
 =\left[B_\infty(1-\alpha)^2c^{1-\alpha}\right]^{1/\alpha},
 \qquad p_X\to\frac1{B_\infty(1-\alpha)}.
 \label{eq:selection-limit}
\end{equation}
Thus distant local maxima do not create an ambiguity in the late global
choice in this region. The result does not assert uniqueness at every state.
\end{proposition}
\begin{proof}
The complementary allocation converges to a finite positive quantity per
unit of $H$ as the common marginal composite product decreases to $c$.
Its inverse demand is smooth there. The substituting allocation dominates
total services. Expanding its CES for large AI supply and using the smooth
complementary allocation gives
\begin{align*}
 g(x)&=cx+O(x^{1-\nu}),&g'(x)&=c+O(x^{-\nu}),\\
 g''(x)&=O(x^{-1-\nu}),&g'''(x)&=O(x^{-2-\nu}).
\end{align*}
For the derivative bounds, one can also use the allocation parameter
$\mu=g'(x)$. The substituting demand has
$x_S\asymp(\mu-c)^{-1/\nu}$ and a differentiable power expansion,
whereas $x_C(\mu)$ and its derivatives are bounded near $c$.
Inverting their sum yields the displayed derivative orders.
Consequently, on every compact positive interval of $y$,
$f_k$ and its first two derivatives converge uniformly to those of
\begin{equation}
 f_\infty(y;B)=(1-\alpha)c^{1-\alpha}y^{1-\alpha}-y/B.
 \label{eq:selection-limit-objective}
\end{equation}
This function is strictly concave, has a positive maximum, and has the
unique maximizer $y_*(B)$ in~\eqref{eq:selection-limit}.

I must also rule out maximizers escaping the compact intervals on which
this convergence holds. Concavity gives $xg'(x)\leq g(x)$ and a supporting
affine bound $g(x)\leq a+bx$ with finite positive constants. Hence
\[
 f_k(y;B)\leq(1-\alpha)(a/k+by)^{1-\alpha}-y/B.
\]
This bound excludes uniformly large $y$, and shows that supplies with
$y\to0$ yield limiting profit per unit of capital at most zero.
The positive value attainable at $y_*(B)$ therefore confines all global
maximizers to a common compact positive interval. Uniform convergence
there and the strict limit maximum imply that all maximizers approach
$y_*(B)$. On a small neighborhood of that limit, the second derivative
is strictly negative for large $k$. All global maximizers lie in that
neighborhood, so there is exactly one. These arguments hold uniformly
for $B$ near $B_\infty$. Substitution into inverse demand gives the price limit.
\end{proof}

The smooth first-order equation has a strictly negative second derivative
at this optimum. Thus the implicit function theorem gives a smooth local
supply branch in the high-capital region. The operating-profit envelope
$\pi_B=X^*/B^2$ is valid there. This restores a smooth dated profit function
for the research problem near this limit, but does not prove global
intertemporal optimality of a proposed research trajectory.

\section{Conditional growth and wages under monopoly}
Suppose an equilibrium satisfies $K/(AN)\to\infty$, $B\to B_\infty>0$,
and $C/K\to h>0$. Set
\[
 R(B)=\alpha[B(1-\alpha)^2c]^{(1-\alpha)/\alpha}-\delta.
\]
The preceding proposition gives $Y/K\to R(B_\infty)/\alpha+\delta/\alpha$
and $r\to R(B_\infty)$. The household Euler equation then implies limiting
average output-per-worker growth $g=R(B_\infty)-\rho$.
For the strictly accelerated regime require $g>\gamma$.

Let $Q=(1-\alpha)(K/(Z_S+Z_C))^\alpha$ and
$p_*=1/[B_\infty(1-\alpha)]$. Then $Q\to Q_*=p_*/c$.
The complementary activity has $X_C/(AL_C)$ converging to a finite positive
constant, since its marginal composite product at zero AI exceeds one
and $p_*/Q_*=c<1$. Consequently,
\begin{equation}
 \frac{w_C}{A}\to
 \left[\frac{Q_*^{1-\sigma_C}-\omega_{XC}^{\sigma_C}p_*^{1-\sigma_C}}
 {\omega_{LC}^{\sigma_C}}\right]^{1/(1-\sigma_C)}>0.
 \label{eq:selection-C-wage}
\end{equation}
To obtain this expression, solve the CES unit-cost equality for the price
of effective labor. In $S$, differentiate production directly:
\[
 w_S=AQ\omega_{LS}\left(\frac{Z_S}{AL_S}\right)^{1/\sigma_S}.
\]
Since $Z_C/(AL_C)$ stays finite and $Z_S\sim cX_S$ grows at the aggregate
scale, the limiting average growth rates satisfy
\begin{equation}
 \boxed{g_{w_C}=\gamma,\qquad
 g_{w_S}=\gamma+\frac{g-\gamma}{\sigma_S}.}
 \label{eq:selection-wages}
\end{equation}
Thus $w_S/w_C\to\infty$, $Y_C/Y\to0$, and both labor-income shares in
aggregate output vanish. The complementary activity remains active.
These are not welfare comparisons or claims about the sign of transitional
wage changes. Instantaneous-rate convergence requires further regularity.

\section{Can efficiency stop below its upper bound?}
The answer is no within a regular accelerated equilibrium with a smooth
normal developer optimum. I state the extra regularity explicitly.

\begin{proposition}[Attainment within a regular accelerated equilibrium]
Suppose an equilibrium has $K/(AN)\to\infty$, $B\to B_\infty>0$,
$C/K\to h>0$, finite convergent $\dot K/K$ and $r$, and
$\overline B<\infty$. Suppose the developer has a normal differentiable
optimum with a nonnegative current-value costate satisfying its necessary
research and costate conditions at sufficiently late dates.
Then $B_\infty=\overline B$.
\end{proposition}
\begin{proof}
Assume instead $B_\infty<\overline B$. Static uniqueness at late dates
gives $\pi_B=X^*/B^2$, with $X^*/K\to y_*(B_\infty)>0$.
Write $F(B,M)=\chi(BM)^\eta\psi(B)$ and
$\psi(B)=1-B/\overline B$. The costate equation is
\[
 \dot q=(r-F_B)q-\pi_B,\qquad
 F_B=\dot B\left(\eta/B+\psi'/\psi\right).
\]
Integrate this linear equation from $t$ to $t+1$. Since $q_{t+1}\geq0$,
\[
 q_t\geq\int_t^{t+1}
 \exp\!\left[-\int_t^s(r_v-F_{B,v})\,dv\right]\pi_{B,s}\,ds.
\]
Because $B$ converges monotonically below its upper bound, the coefficient
$\eta/B+\psi'/\psi$ is bounded, and the integral of $|F_B|$ on
$[t,t+1]$ tends to zero. Convergence of $r$ and $\dot K/K$ bounds the
other factors on this interval. Thus $q_t\geq aK_t$ for some $a>0$
at all sufficiently late dates. This lower bound makes research interior,
since $0<\eta<1$ and the marginal research benefit at $M=0$ is infinite.
The necessary condition gives
\[
 M_t=[q_t\chi\eta B_t^\eta\psi(B_t)]^{1/(1-\eta)}
 \geq bK_t^{1/(1-\eta)}
\]
for some $b>0$. The resource constraint, the finite $Y/K$ limit, positivity
of consumption and compute use, and bounded $\dot K/K$ instead imply
$M_t/K_t$ is bounded above. These inequalities contradict $K_t\to\infty$.
Here $K\to\infty$ follows from $K/(AN)\to\infty$, $\gamma>0$ and $n\geq0$.
\end{proof}

This proposition establishes a necessary property of the stated regular
optimal equilibrium. It does not construct one, justify global optimality
from first-order conditions, or establish the required costate regularity.
Under its assumptions the preceding growth and wage formulas may use
$B_\infty=\overline B$. Proving existence and verifying the developer's
global optimum are the next tasks. The mobile-labor case remains separate.
\end{document}
